% Berkas induk kompilasi skripsi sarjana UNESA.
% Kompilasi via terminal: latexmk -pdf Skripsi.tex

\PassOptionsToPackage{table}{xcolor}
\documentclass{SkripsiUnesa}
\usepackage{booktabs}
\usepackage{float}
\usepackage[section]{placeins}
\pgfplotsset{compat=1.18}
\loadglsentries{Istilah.tex}
\raggedbottom

% Aktifkan baris ini untuk mempercepat kompilasi hanya pada bab yang sedang ditulis.
% \includeonly{src/01-body/04-bab4}

\begin{document}

% Memuat metadata skripsi dari DataSkripsi.tex.
% Data dan metadata dokumen skripsi atau proposal UNESA.
% Seluruh data administratif dikelola di sini agar tidak perlu mengubah berkas .cls.

% Judul naskah dalam bahasa Indonesia dan bahasa Inggris.
% Judul bahasa Inggris digunakan untuk sampul dan halaman abstrak bahasa Inggris.
\Judul{Pengembangan \textit{Knowledge Graph-Based Retrieval Augmented Generation} dengan \textit{Hybrid Vector-Graph Retrieval} untuk Asisten Pencarian literatur Akademik}

\JudulEng{Development of \textit{Knowledge Graph-Based Retrieval Augmented Generation} with \textit{Hybrid Vector-Graph Retrieval} for Academic Literature Search Assistant}

% Identitas penulis.
\Nama{Rizky Yanuar Kristianto}
\NIM{22031554017}

% Program studi, tahun kelulusan atau sidang, dan fakultas.
% Fakultas membutuhkan dua argumen: nama lengkap dan singkatan resmi.
\ProgramStudi{Sarjana Sains Data}
\Programme{Bachelor of Data Science}
\Tahun{2026}

\Fakultas{Matematika dan Ilmu Pengetahuan Alam}{MIPA}
\Faculty{Mathematics and Natural Sciences}
\Universitas{Universitas Negeri Surabaya}
\Institution{State University of Surabaya}

% Pejabat program studi dan fakultas.
\KPS{Yuliani Puji Astuti, M.Si.}
\NIPkps{197807312006042001}

\Dekan{Prof. Dr. Wasis, M.Si.}
\NIPDekan{196712031993021001}

% Dosen pembimbing dan NIP.
% Format: {Pembimbing 1}{Pembimbing 2}. Kosongkan kurung kedua {} jika hanya satu pembimbing.
\Pembimbing{Ibnu Febry K., S.Kom., M.Sc., Ph.D.}
{Hasanuddin Al-Habib, M.Si.}

\NIPPembimbing{198802182015041001}
{199308092022031010}

% Dosen penguji dan NIP.
% Format: {Ketua Penguji}{Anggota 1}{Anggota 2}. Kosongkan kurung ketiga {} jika hanya dua penguji.
\Penguji{Dr. Atik Wintarti, M.Kom.}
{Yuni Rosita Dewi, M.Si.}
{}

\NIPPenguji{196610121991032002}
{199506282024062001}
{}

% Pembimbing ahli atau mitra luar instansi (isi tanda strip jika tidak ada).
\Ahli{-}
\NIPahli{-}

% Tanggal persetujuan, ujian sidang, dan pernyataan orisinalitas.
\TanggalUji{3 juli 2026}
\TanggalDisetujui{17 juni 2026}
\TanggalSidang{3 Juli 2026}
\TanggalOrisinalitas{17 Juli 2026}

% Deklarasi penggunaan kecerdasan buatan pada surat pernyataan AI.
% Gunakan \checkedbox untuk mencentang atau \emptybox untuk membiarkan kotak kosong.
% Jika mencentang opsi lainnya, isi rincian penggunaannya pada \AILainnyaTeks.
\renewcommand{\AIEksplorasiIde}{\checkedbox}
\renewcommand{\AIKerangkaTulisan}{\checkedbox}
\renewcommand{\AITataBahasa}{\emptybox}
\renewcommand{\AIAnalisisData}{\checkedbox}
\renewcommand{\AIIlustrasi}{\emptybox}
\renewcommand{\AILainnya}{\emptybox}
\renewcommand{\AILainnyaTeks}{}


% Bagian awal: penomoran halaman romawi kecil (i, ii, iii, ...).
\Awal

% Halaman administratif sesuai urutan resmi panduan skripsi UNESA.
\HalamanJudul
\SampulDalam
\HalamanPersetujuan
\HalamanPengesahan
\SuratOrisinalitas
\SuratPernyataanAI

% Abstrak bahasa Indonesia, bahasa Inggris, dan prakata.
\begin{Abstrak}
    Produksi literatur ilmiah di lingkungan perguruan tinggi menunjukkan peningkatan yang signifikan setiap tahunnya. Peningkatan jumlah publikasi menyebabkan peneliti mengalami kesulitan dalam menemukan literatur yang relevan secara efisien. Sistem pencarian akademik konvensional yang mengandalkan pencocokan kata kunci kurang mampu memahami konteks semantik dan relasi antar-entitas akademik. Penelitian ini bertujuan merancang \textit{pipeline} konstruksi \textit{Knowledge Graph} dari korpus literatur akademik, merancang mekanisme \textit{Hybrid Vector-Graph Retrieval} dalam arsitektur \textit{Knowledge Graph-Based Retrieval Augmented Generation}, dan mengevaluasi efektivitas integrasi \textit{Knowledge Graph} dalam meningkatkan kualitas jawaban model bahasa besar. Data publikasi diperoleh dari profil dosen dan publikasi rumpun Informatika dan Komputer UNESA. Konstruksi graf menghasilkan 45.685 simpul dan 122.675 relasi dengan 16 jenis simpul yang mencakup entitas bibliografis dan konseptual. Eksperimen dilakukan dengan membandingkan tiga metode \textit{retrieval}, yaitu \textit{Naive}, \textit{Subgraph}, dan \textit{Hybrid} pada 90 pertanyaan evaluasi yang terbagi ke dalam kategori faktual, relasional, dan \textit{multi-hop}. Evaluasi \textit{retrieval} menggunakan \textit{Mean Reciprocal Rank} dan \textit{Hit@K}, sedangkan evaluasi generasi jawaban menggunakan metode \textit{Pairwise LLM as Judge}. Metode \textit{Hybrid} menghasilkan performa terbaik dengan \textit{MRR} sebesar 0,50 dan \textit{Win Rate Overall} sebesar 53,9\% terhadap metode \textit{Naive}. Metode \textit{Subgraph} memperoleh \textit{MRR} sebesar 0,43 dengan \textit{Win Rate Overall} sebesar 45,3\% terhadap \textit{Naive}. Evaluasi pada dimensi \textit{Traceability} dan \textit{Faithfulness} menunjukkan integrasi \textit{Knowledge Graph} membantu model bahasa besar menghasilkan jawaban yang lebih mudah dilacak sumbernya dan mengurangi klaim yang tidak terverifikasi.

    \katakunci{\textit{Knowledge Graph}, \textit{Retrieval Augmented Generation}, \textit{Hybrid Retrieval}, Pencarian Literatur Akademik, Model Bahasa Besar}
\end{Abstrak}

\begin{Abstract}
    Scientific literature production in higher education institutions shows a significant increase every year. The growing number of publications causes researchers to experience difficulties in finding relevant literature efficiently. Conventional academic search systems that rely on keyword matching are less capable of understanding semantic context and relationships between academic entities. This study aims to design a Knowledge Graph construction pipeline from academic literature corpora, design a Hybrid Vector-Graph Retrieval mechanism within a Knowledge Graph-Based Retrieval Augmented Generation architecture, and evaluate the effectiveness of Knowledge Graph integration in improving the quality of large language model responses. Publication data were obtained from lecturer profiles and publications of the Informatics and Computer Science department at UNESA. The graph construction produced 45,685 nodes and 122,675 relationships with 16 node types covering bibliographic and conceptual entities. Experiments were conducted by comparing three retrieval methods, namely Naive, Subgraph, and Hybrid on 90 evaluation questions divided into factual, relational, and multi-hop categories. Retrieval evaluation used Mean Reciprocal Rank and Hit@K, while answer generation evaluation used the Pairwise LLM as Judge method. The Hybrid method achieved the best performance with an MRR of 0.50 and an Overall Win Rate of 53.9\% against the Naive method. The Subgraph method obtained an MRR of 0.43 with an Overall Win Rate of 45.3\% against Naive. Evaluation on the Traceability and Faithfulness dimensions showed that Knowledge Graph integration helps large language models produce answers that are more traceable to their sources and reduces unverified claims.

    \keywords{Knowledge Graph, Retrieval Augmented Generation, Hybrid Retrieval, Academic Literature Search, Large Language Model}
\end{Abstract}

\Prakata

Puji syukur kehadirat Allah Swt. atas segala rahmat, nikmat, dan pertolongan-Nya sehingga penulis dapat menyelesaikan skripsi ini. Dalam proses penyusunannya, penulis menyadari bahwa banyak pihak yang telah memberikan dukungan, bantuan, dan semangat. Oleh karena itu, penulis menyampaikan terima kasih kepada:

\begin{enumerate}[label=\arabic*., leftmargin=0.75cm, itemsep=6pt]
    \item Bapak Ibnu Febry K., S.Kom., M.Sc., Ph.D. dan Bapak Hasanuddin Al-Habib, M.Si. selaku dosen pembimbing yang telah memberikan arahan, bimbingan, masukan, dan kesabaran selama proses penyusunan skripsi ini.
    \item Ibu Dr. Atik Wintarti, M.Kom. dan Ibu Yuni Rosita Dewi, M.Si. selaku dosen penguji yang telah memberikan kritik, saran, dan masukan yang membangun untuk penyempurnaan skripsi ini.
    \item Seluruh teman angkatan 2022 Program Studi Sarjana Sains Data Universitas Negeri Surabaya yang tidak dapat disebutkan satu per satu, terima kasih atas dukungan, bantuan, dan kebersamaan selama masa perkuliahan.
\end{enumerate}

Akhir kata, penulis berharap skripsi ini dapat memberikan manfaat bagi perkembangan ilmu pengetahuan, khususnya di bidang Sains Data.

\vspace{1cm}
\begin{flushright}
    Surabaya, Juni 2026\\
    \vspace{1.5cm}
    Penulis
\end{flushright}


% Daftar isi, tabel, gambar, dan simbol.
\DaftarIsi
\DaftarTabel
\DaftarGambar
% Daftar simbol matematika dan notasi naskah.
% Tambahkan baris baru dengan format: $simbol$ & penjelasan simbol \\.
\DaftarSimbol
\begin{longtable}{p{0.15\textwidth} p{0.75\textwidth}}
    % Struktur Pengetahuan
    $G$                                     & \textit{Knowledge Graph} akademik, terdiri atas himpunan entitas ($E$), relasi ($R$), dan \textit{triplet} ($T$).   \\
    $(h, r, t)$                             & \textit{Triplet} pengetahuan: \textit{head} (subjek), \textit{relation} (predikat), \textit{tail} (objek).          \\
    $\mathcal{S}$                           & Himpunan \textit{triplet} hasil pemetaan fungsi ekstraksi informasi $f_{IE}$ dari teks $\mathcal{T}$.               \\
    $E$                                     & Himpunan entitas dalam \textit{Knowledge Graph}.                                                                    \\
    $R$                                     & Himpunan relasi dalam \textit{Knowledge Graph}.                                                                     \\

    % Pemodelan Academic RAG
    $q$ / $Q$                               & Pertanyaan atau kueri pengguna.                                                                                     \\
    $K^{\text{content}}$                    & Himpunan kata kunci korpus akademik.                                                                                \\
    $K^{\text{clues}}$                      & Kandidat \textit{clues} hasil penyaringan kata kunci berdasarkan ambang kemiripan $\tau$.                           \\
    $f_{\text{embed}}(\cdot)$               & Fungsi \textit{embedding} yang memetakan teks ke representasi vektor.                                               \\
    $\tau$                                  & Ambang batas (\textit{threshold}) kemiripan untuk penyaringan \textit{clues}.                                       \\
    $\mathcal{K}_l, \mathcal{K}_h$          & Kata kunci hierarkis; $\mathcal{K}_l$ untuk pencarian lokal detail, $\mathcal{K}_h$ untuk konteks global.           \\
    $\mathcal{L}_{\text{local}}$            & Konteks lokal berupa entitas, relasi, dan unit teks dari subgraf.                                                   \\
    $\mathcal{L}_{\text{global}}$           & Konteks global berupa entitas, relasi, dan unit teks dari jejaring lintas domain.                                   \\
    $\mathcal{E}, \mathcal{R}, \mathcal{T}$ & Komponen konteks: entitas ($\mathcal{E}$), relasi ($\mathcal{R}$), unit teks ($\mathcal{T}$).                       \\

    % LLM
    $P_\theta(y \mid y_{<t}, x)$            & Distribusi probabilitas bersyarat token pada model bahasa besar dengan parameter $\theta$.                          \\

    % Metrik Evaluasi
    $|Q|$                                   & Jumlah kueri dalam himpunan evaluasi.                                                                               \\
    $\text{rank}_i$                         & Posisi dokumen relevan pertama pada kueri ke-$i$.                                                                   \\
    $Q_{\text{relevan@}K}$                  & Himpunan kueri yang memiliki dokumen relevan dalam $K$ peringkat teratas.                                           \\
    $|V|$                                   & Jumlah klaim terverifikasi dalam jawaban, untuk menghitung \textit{Faithfulness}.                                   \\
    $|S|$                                   & Total klaim yang diekstraksi dari jawaban sistem.                                                                   \\
    $E_q, E_{q_i}$                          & Vektor \textit{embedding} kueri asli ($E_q$) dan kueri sintetis ke-$i$ ($E_{q_i}$) untuk \textit{Answer Relevancy}. \\
    $m$                                     & Jumlah pertanyaan sintetis yang dibangkitkan untuk menghitung \textit{Answer Relevancy}.                            \\
    $M$                                     & Metode pencarian yang dievaluasi (misal: \textit{Naive}, \textit{Subgraph}, \textit{Hybrid}).                       \\
\end{longtable}


% Bagian inti: penomoran halaman arab dimulai dari angka 1.
\Inti

% Bab 1 sampai Bab 5.
\include{src/01-body/01-bab1}
\include{src/01-body/02-bab2}
\include{src/01-body/03-bab3}
\include{src/01-body/04-bab4}
\include{src/01-body/05-bab5}

% Bagian akhir: daftar pustaka (dari Pustaka.bib) dan lampiran.
\DaftarPustaka{Pustaka}
\BukaLampiran

\lampiran{\textit{Source code}}

\lstset{
    basicstyle=\ttfamily\small,
    keywordstyle=\bfseries,
    commentstyle=\itshape\color{gray},
    breaklines=true,
    frame=single,
    numbers=left,
    numberstyle=\tiny,
    xleftmargin=2em,
    framexleftmargin=1.5em,
    tabsize=4,
}

\begin{lstlisting}[language=Python, caption={Algoritma \textit{Hybrid Vector-Graph Retrieval}}]
ALGORITHM: HybridVectorGraphRetrieval

INPUT:
    Q           : kueri pengguna
    VectorDB    : basis data vektor (Zilliz Cloud)
    GraphDB     : basis data graf (Neo4j AuraDB)
    Reranker    : model Cross-Encoder Reranker
    mode        : "naive" | "subgraph" | "hybrid"
    top_k       : jumlah dokumen yang diambil

OUTPUT:
    context     : konteks terstruktur untuk LLM

BEGIN
    result = {}

    # === TAHAP 1: Dekomposisi Kata Kunci ===
    IF mode IN {subgraph, hybrid} THEN
        keyword_clues = VectorDB.search("ContentKeyword", Q)
        (K_high, K_low) = LLM_decompose(Q, keyword_clues)
        local_query  = combine(K_low, Q)
        global_query = combine(K_high, Q)
    END IF

    # === TAHAP 2: Pencarian Vektor (Naive) ===
    IF mode IN {naive, hybrid} THEN
        paper_chunks = VectorDB.search("PaperChunk", Q, top_k * 3)
    END IF

    # === TAHAP 3: Pencarian Subgraf Lokal ===
    IF mode IN {subgraph, hybrid} THEN
        entities = EntitySubgraphRetrieval(local_query, VectorDB,
                                           GraphDB, top_k)
        subgraph_triples = GraphDB.shortest_path(entities, depth=2)
        text_units = VectorDB.search("PaperChunk", local_query)
    END IF

    # === TAHAP 4: Pencarian Relasi Global ===
    IF mode IN {hybrid} THEN
        relationships = VectorDB.search("RelationshipEmbedding",
                                         global_query, top_k * 2)
        FOR EACH rel IN relationships DO
            src_deg = GraphDB.get_degree(rel.src_id)
            tgt_deg = GraphDB.get_degree(rel.tgt_id)
            rel.distance = rel.distance /
                (1 + 0.05 * (log(src_deg+1) + log(tgt_deg+1)))
        END FOR
        relationships = sort_ascending(relationships, key=distance)
    END IF

    # === TAHAP 5: Fusi dan Pemeringkatan Ulang ===
    IF mode == hybrid THEN
        ranked_papers = RRF_Rerank(paper_chunks, entities,
                                    GraphDB, Reranker, Q, top_k)
    ELSE IF mode == naive THEN
        ranked_papers = Reranker.rerank(Q, paper_chunks, top_k)
    END IF

    # === TAHAP 6: Penyusunan Konteks ===
    context = format_context(
        paper_chunks  = ranked_papers,
        entities      = entities,
        relationships = relationships,
        subgraph      = subgraph_triples,
        text_units    = text_units
    )

    RETURN context
END
\end{lstlisting}


\end{document}
